% GENERATED FILE. Edit canonical lessons, then run npm run generate:books.
% canonical-source-hash: fnv1a64:e055d2f4c933e4b2
% canonical-lessons: TA-C413-vimani, TA-C413-taraiyiranku, TA-C413-irakkai, TA-C413-cunkam, TA-C413-kutivaravu

\chapter{Pilots and Customs}
\label{ch:ta-pilots-and-customs}
\hlchaptermodality{413}

\begin{chapteropening}
\textbf{By the end of this chapter:} \emph{I can talk about a pilot, landing, a wing, customs and immigration.}

Complete the last lesson of chapter 413: I can talk about a pilot, landing, a wing, customs and immigration.
\end{chapteropening}

\section[\texorpdfstring{\ta{விமானி} (vimāṉi)}{விமானி (vimāṉi)}]{\ta{விமானி} (vimāṉi) — a pilot}

\label{lesson:TA-C413-vimani}



\begin{quote}
Before the new one: say the Tamil for to move forward, then the Tamil for to move away.
\end{quote}

\subsection*{You'll want to know: \ta{விமானி}}
\textbf{\ta{விமானி}} — \emph{vimāṉi} — \textquotedblleft{}a pilot\textquotedblright{}.

\begin{quote}
\ta{விமானி} \ta{பயணிகளுக்கு} \ta{வணக்கம்} \ta{சொன்னார்}.
\end{quote}

\begin{quote}
\emph{vimāṉi payaṇikaḷukku vaṇakkam coṉṉār.}
\end{quote}

\begin{quote}
The pilot greeted the passengers.
\end{quote}

\ta{விமானம்}, a plane, + -\ta{இ}, one who: the one who flies it.

\begin{cousinweb}
\textbf{\ta{விமானி}} (\emph{vimāṉi}) is built from \ta{விமானம்}, \textquotedblleft{}an aeroplane\textquotedblright{}, borrowed from Sanskrit \emph{vimāna}, \textquotedblleft{}a car of the gods\textquotedblright{} (CDIAL 11823), with the ending \emph{-i}, \textquotedblleft{}one who\textquotedblright{}.
\end{cousinweb}

\begin{grammarlens}[title={What you've built}]
One more word for getting around.
\end{grammarlens}

\subsection*{Your turn}
\begin{itemize}
\raggedright
  \item \emph{Say it:} \emph{vimāṉi}
  \item \emph{Say it:} \emph{vimāṉi}, once more
  \item \emph{Say it:} \emph{vimāṉi payaṇikaḷukku vaṇakkam coṉṉār.} — The pilot greeted the passengers.
  \item \emph{Recall:} say the Tamil for to move forward, then the Tamil for to move away, then say \emph{vimāṉi} again
\end{itemize}

\begin{tcolorbox}[breakable,colback=teal!4,colframe=teal!35!black,title={Before you move on}]
What does \ta{விமானி} mean? (\textquotedblleft{}A pilot\textquotedblright{}.) Say it once more.
\end{tcolorbox}
\section[\texorpdfstring{\ta{தரையிறங்கு} (taraiyiṟaṅku)}{தரையிறங்கு (taraiyiṟaṅku)}]{\ta{தரையிறங்கு} (taraiyiṟaṅku) — to make a landing; to land}

\label{lesson:TA-C413-taraiyiranku}



\begin{quote}
Before the new one: say the Tamil for to move away, then the Tamil for a pilot.
\end{quote}

\subsection*{You'll want to know: \ta{தரையிறங்கு}}
\textbf{\ta{தரையிறங்கு}} — \emph{taraiyiṟaṅku} — \textquotedblleft{}to make a landing; to land\textquotedblright{}.

\begin{quote}
\ta{விமானம்} \ta{சரியான} \ta{நேரத்தில்} \ta{தரையிறங்கியது}.
\end{quote}

\begin{quote}
\emph{vimāṉam cariyāṉa nērattil taraiyiṟaṅkiyatu.}
\end{quote}

\begin{quote}
The plane landed on time.
\end{quote}

Past (it …): \ta{தரையிறங்கியது}, \emph{taraiyiṟaṅkiyatu}.

\ta{தரை}, the ground, + \ta{இறங்கு}, to come down.

\begin{cousinweb}
\textbf{\ta{தரையிறங்கு}} (\emph{taraiyiṟaṅku}) is built from \ta{தரை}, \textquotedblleft{}the ground, the floor\textquotedblright{}, and \ta{இறங்கு}, \textquotedblleft{}to descend\textquotedblright{} (DEDR 516).
\end{cousinweb}

\begin{grammarlens}[title={What you've built}]
One more word for getting around.
\end{grammarlens}

\subsection*{Your turn}
\begin{itemize}
\raggedright
  \item \emph{Say it:} \emph{taraiyiṟaṅku}
  \item \emph{Say it:} \emph{taraiyiṟaṅku}, once more
  \item \emph{Say it:} \emph{vimāṉam cariyāṉa nērattil taraiyiṟaṅkiyatu.} — The plane landed on time.
  \item \emph{Recall:} say the Tamil for to move away, then the Tamil for a pilot, then say \emph{taraiyiṟaṅku} again
\end{itemize}

\begin{tcolorbox}[breakable,colback=teal!4,colframe=teal!35!black,title={Before you move on}]
What does \ta{தரையிறங்கு} mean? (\textquotedblleft{}To make a landing; to land\textquotedblright{}.) Say it once more.
\end{tcolorbox}
\section[\texorpdfstring{\ta{இறக்கை} (iṟakkai)}{இறக்கை (iṟakkai)}]{\ta{இறக்கை} (iṟakkai) — a wing}

\label{lesson:TA-C413-irakkai}



\begin{quote}
Before the new one: say the Tamil for a pilot, then the Tamil for to make a landing.
\end{quote}

\subsection*{You'll want to know: \ta{இறக்கை}}
\textbf{\ta{இறக்கை}} — \emph{iṟakkai} — \textquotedblleft{}a wing\textquotedblright{}.

\begin{quote}
\ta{சன்னலிலிருந்து} \ta{விமானத்தின்} \ta{இறக்கை} \ta{தெரிந்தது}.
\end{quote}

\begin{quote}
\emph{caṉṉaliliruntu vimāṉattiṉ iṟakkai terintatu.}
\end{quote}

\begin{quote}
The plane's wing could be seen from the window.
\end{quote}

A bird has \ta{இறக்கைகள்}, wings; one feather is \ta{இறகு}.

\begin{cousinweb}
\textbf{\ta{இறக்கை}} (\emph{iṟakkai}) is native Dravidian (DEDR 2591), in the same article as \ta{இறகு}, \textquotedblleft{}a feather\textquotedblright{}.
\end{cousinweb}

\begin{grammarlens}[title={What you've built}]
One more word for getting around.
\end{grammarlens}

\subsection*{Your turn}
\begin{itemize}
\raggedright
  \item \emph{Say it:} \emph{iṟakkai}
  \item \emph{Say it:} \emph{iṟakkai}, once more
  \item \emph{Say it:} \emph{caṉṉaliliruntu vimāṉattiṉ iṟakkai terintatu.} — The plane's wing could be seen from the window.
  \item \emph{Recall:} say the Tamil for a pilot, then the Tamil for to make a landing, then say \emph{iṟakkai} again
\end{itemize}

\begin{tcolorbox}[breakable,colback=teal!4,colframe=teal!35!black,title={Before you move on}]
What does \ta{இறக்கை} mean? (\textquotedblleft{}A wing\textquotedblright{}.) Say it once more.
\end{tcolorbox}
\section[\texorpdfstring{\ta{சுங்கம்} (cuṅkam)}{சுங்கம் (cuṅkam)}]{\ta{சுங்கம்} (cuṅkam) — customs; a toll}

\label{lesson:TA-C413-cunkam}



\begin{quote}
Before the new one: say the Tamil for to make a landing, then the Tamil for a wing.
\end{quote}

\subsection*{You'll want to know: \ta{சுங்கம்}}
\textbf{\ta{சுங்கம்}} — \emph{cuṅkam} — \textquotedblleft{}customs; a toll\textquotedblright{}.

\begin{quote}
\ta{விமான} \ta{நிலையத்தில்} \ta{சுங்க} \ta{அதிகாரிகள்} \ta{பைகளைச்} \ta{சோதித்தார்கள்}.
\end{quote}

\begin{quote}
\emph{vimāṉa nilaiyattil cuṅka atikārikaḷ paikaḷaic cōtittārkaḷ.}
\end{quote}

\begin{quote}
At the airport the customs officers checked the bags.
\end{quote}

As a modifier it drops -\ta{ம்}: \ta{சுங்க} \ta{அதிகாரி}, a customs officer.

\begin{cousinweb}
\textbf{\ta{சுங்கம்}} (\emph{cuṅkam}) is borrowed from Prakrit \emph{suṅka}, \textquotedblleft{}a toll, a duty\textquotedblright{}, from Sanskrit \emph{śulka}, \textquotedblleft{}a price, a tax\textquotedblright{} (CDIAL 12543).
\end{cousinweb}

\begin{grammarlens}[title={What you've built}]
One more word for getting around.
\end{grammarlens}

\subsection*{Your turn}
\begin{itemize}
\raggedright
  \item \emph{Say it:} \emph{cuṅkam}
  \item \emph{Say it:} \emph{cuṅkam}, once more
  \item \emph{Say it:} \emph{vimāṉa nilaiyattil cuṅka atikārikaḷ paikaḷaic cōtittārkaḷ.} — At the airport the customs officers checked the bags.
  \item \emph{Recall:} say the Tamil for to make a landing, then the Tamil for a wing, then say \emph{cuṅkam} again
\end{itemize}

\begin{tcolorbox}[breakable,colback=teal!4,colframe=teal!35!black,title={Before you move on}]
What does \ta{சுங்கம்} mean? (\textquotedblleft{}Customs; a toll\textquotedblright{}.) Say it once more.
\end{tcolorbox}
\section[\texorpdfstring{\ta{குடிவரவு} (kuṭivaravu)}{குடிவரவு (kuṭivaravu)}]{\ta{குடிவரவு} (kuṭivaravu) — immigration}

\label{lesson:TA-C413-kutivaravu}



\begin{quote}
Before the new one: say the Tamil for a wing, then the Tamil for customs.
\end{quote}

\subsection*{You'll want to know: \ta{குடிவரவு}}
\textbf{\ta{குடிவரவு}} — \emph{kuṭivaravu} — \textquotedblleft{}immigration\textquotedblright{}.

\begin{quote}
\ta{குடிவரவு} \ta{வரிசையில்} \ta{ஒரு} \ta{மணி} \ta{நேரம்} \ta{நின்றோம்}.
\end{quote}

\begin{quote}
\emph{kuṭivaravu varicaiyil oru maṇi nēram niṉṟōm.}
\end{quote}

\begin{quote}
We stood in the immigration queue for an hour.
\end{quote}

\ta{குடி}, a home, + \ta{வரவு}, coming: coming to live.

\begin{cousinweb}
\textbf{\ta{குடிவரவு}} (\emph{kuṭivaravu}) is built from \ta{குடி}, \textquotedblleft{}a house, a home\textquotedblright{} (DEDR 1655), and \ta{வரவு}, \textquotedblleft{}coming\textquotedblright{} (DEDR 5270).
\end{cousinweb}

\begin{grammarlens}[title={What you've built}]
One more word for getting around.
\end{grammarlens}

\subsection*{Your turn}
\begin{itemize}
\raggedright
  \item \emph{Say it:} \emph{kuṭivaravu}
  \item \emph{Say it:} \emph{kuṭivaravu}, once more
  \item \emph{Say it:} \emph{kuṭivaravu varicaiyil oru maṇi nēram niṉṟōm.} — We stood in the immigration queue for an hour.
  \item \emph{Recall:} say the Tamil for a wing, then the Tamil for customs, then say \emph{kuṭivaravu} again
\end{itemize}

\begin{tcolorbox}[breakable,colback=teal!4,colframe=teal!35!black,title={Before you move on}]
What does \ta{குடிவரவு} mean? (\textquotedblleft{}Immigration\textquotedblright{}.) Say it once more.
\end{tcolorbox}
